\subsubsection{Perceptrón multicapa - MLP}

Un perceptrón multicapa es una red neuronal que se utiliza en aprendizaje supervisado en el campo del aprendizaje automático. El MLP está compuesto por múltiples capas de neuronas interconectadas, en las que las salidas de las neuronas de una capa se convierten en entradas para la siguiente capa. La primera capa se llama capa de entrada, la última capa se llama capa de salida y las capas intermedias se llaman capas ocultas. El MLP presenta una gran capacidad para modelar relaciones no lineales y para generalizar. Como punto negativo del uso de este modelo, podemos remarcar la necesidad de un gran conjunto de datos de entrenamiento para evitar el sobreajuste. Por último, cabe resaltar que los MLP's fueron los precursores de otras redes neuronales más complejas como la redes convolucionales y las redes recurrentes. (Fuente: https://gamco.es/glosario/percepatron-multicapa-mlp/) %Hay que poner la fuente en el .bib